\subsection*{Цель работы}
Практическое ознакомление с пассивными и активными дифференцирующими и интегрирующими цепями как схемами формирования сигналов и частотными фильтрами.

\subsection*{Теоретические сведения}

Дифференцирующие (ДЦ) и интегрирующие (ИЦ) цепи применяются для изменения формы импульсных сигналов. 

\begin{enumerate}
    \item \textit{Пассивная ДЦ.} Представляет собой $RC$-цепь, где выходной сигнал снимается с резистора $R$ (\cref{fig:passive_dc}). Постоянная времени цепи $\tau = RC$. Для качественного дифференцирования необходимо выполнение условия $3\tau < \tau_и$, где $\tau_и$ --- длительность импульса. Ошибка дифференцирования определяется как:
    \begin{equation*}
        \varepsilon_д = \frac{3\tau}{\tau_и} \cdot 100\%
    \end{equation*}
    С точки зрения частотного анализа ДЦ является фильтром высоких частот (ФВЧ).

    \item \textit{Пассивная ИЦ.} Выходной сигнал снимается с конденсатора $C$ (\cref{fig:passive_ic}). Условие качественного интегрирования: $3\tau > \tau_и$. Ошибка интегрирования:
    \begin{equation*}
        \varepsilon_и = \frac{\tau_и}{3\tau} \cdot 100\%
    \end{equation*}
    ИЦ является фильтром низких частот (ФНЦ).

    \item \textit{Активные цепи.} Реализуются на базе операционных усилителей (ОУ). Передаточная характеристика определяется соотношением сопротивлений обратной связи $Z_{ос}$ и входной цепи $Z_{вх}$:
    \begin{equation*}
        K(jf) = -\frac{Z_{ос}(jf)}{Z_{вх}(jf)}
    \end{equation*}
    Активные цепи позволяют избежать затухания сигнала и обеспечивают более высокую точность выполнения математических операций.
\end{enumerate}

\subsection*{Инструкция по выполнению работы}

\subsubsection*{Подготовка установки}
\begin{enumerate}
    \item Подайте напряжение питания $\pm 15$~В на макет.
    \item Подключите выход генератора к входу исследуемой цепи, а вход осциллографа --- к ее выходу.
    \item Установите параметры $R$ и $C$ на лицевой панели макета согласно указаниям преподавателя.
\end{enumerate}

\subsubsection*{Снятие АЧХ (п. 2 программы)}
\begin{enumerate}
    \item Установите на генераторе гармонический сигнал (синусоида) амплитудой $1\dots 2$~В.
    \item Изменяя частоту $f$ в диапазоне от $10$~Гц до $1$~МГц, фиксируйте амплитуду выходного напряжения $U_{вых}$ для каждой из четырех цепей (пассивные ДЦ/ИЦ, активные ДЦ/ИЦ).
    \item Данные заносите в \cref{tab:ach_data}.
    \item Для каждой цепи найдите граничную частоту $f_{гр}$, на которой $U_{вых}$ уменьшается до уровня $0.7$ от максимального значения. Запишите в \cref{tab:f_gr}.
\end{enumerate}

\subsubsection*{Исследование формы импульсов (п. 3 программы)}
\begin{enumerate}
    \item Переключите генератор в режим формирования прямоугольных импульсов (меандр).
    \item Для пассивной ДЦ: установите малую длительность импульса $\tau_и$ (случай $3\tau < \tau_и$) и зарисуйте осциллограмму. Затем увеличьте $\tau_и$ и зафиксируйте изменения.
    \item Для пассивной ИЦ: добейтесь режима линейного нарастания напряжения (интегрирования), варьируя $\tau_и$.
    \item \textit{Важно:} при зарисовке фиксируйте амплитудные значения $U_0$ и $U_{Cm}$ для последующего расчета ошибок по формулам из теории.
    \item Повторите измерения для активных цепей. Результаты занесите в \cref{tab:impulse_data}.
\end{enumerate}

\subsection*{Инструкция по обработке результатов}
\begin{enumerate}
    \item По данным \cref{tab:ach_data} постройте графики АЧХ в логарифмическом масштабе по оси частот.
    \item Рассчитайте теоретические значения граничных частот $f_{гр} = 1/(2\pi RC)$ и сравните их с экспериментальными.
    \item Рассчитайте ошибки дифференцирования $\varepsilon_д$ и интегрирования $\varepsilon_и$ на основе зафиксированных параметров импульсов.
\end{enumerate}

\begin{figure}[H]
    \centering
    \includegraphics[width=0.4\linewidth]{figs/scheme_dc.pdf}
    \caption{Схема пассивной дифференцирующей цепи}
    \label{fig:passive_dc}
\end{figure}

\begin{figure}[H]
    \centering
    \includegraphics[width=0.4\linewidth]{figs/passive_ic.pdf}
    \caption{Схема пассивной интегрирующей цепи}
    \label{fig:passive_ic}
\end{figure}

\begin{figure}[H]
    \centering
    \includegraphics[width=0.7\linewidth]{figs/active_circuits.pdf}
    \caption{Схемы активных цепей на ОУ: а) дифференцирующая, б) интегрирующая}
    \label{fig:active_circuits}
\end{figure}